% \section{Implementations of Dataset Generation}
% \label{app:dataset}

% This appendix provides additional details of the dataset-generation
% procedure introduced in Sec.~\ref{sec:dataset}. We describe how the
% original Text-to-CadQuery samples are converted into physics-aware
% inverse-design tasks, including CAD--CAE annotation, automatic physical
% condition construction, design-domain generation, geometry semantic
% construction, and dataset partitioning. We also provide representative
% examples of the final model-visible prompts.

% \subsection{Reference CAD--CAE Annotation}
% \label{app:dataset_annotation}

% We start from the Text-to-CadQuery corpus~\citep{xie2025text}, which
% contains approximately \(110\text{k}\) CadQuery programs paired with
% geometric descriptions. Each source program is treated as a
% reference design: it provides a valid geometric realization from
% which an inverse-design task can be constructed, but its exact geometry
% is not exposed to the model.

% Each CadQuery program is first executed to construct the corresponding
% solid. Samples are discarded if the program cannot be executed, does
% not produce exactly one valid connected solid, or cannot be processed by
% the downstream geometry and simulation pipeline. Successfully
% constructed solids are exported and analyzed using the same deterministic
% finite-element pipeline.

% All samples are evaluated using three-dimensional small-strain linear
% elasticity in SI units. We use isotropic structural steel with Young's
% modulus \(2.1\times10^{11}\,\mathrm{Pa}\), Poisson's ratio \(0.3\), and
% density \(7850\,\mathrm{kg/m^3}\). For each successfully simulated
% reference design, we record the maximum nodal displacement
% \(\delta_{\max}\), the volume-weighted \(99\)-th percentile element
% von Mises stress \(\sigma_{\mathrm{P99}}\), and the total mass \(m\).
% These quantities are used only to construct the performance requirements
% of the inverse-design task and are not directly revealed to the model.


% \subsection{Automatic Physical Conditions}
% \label{app:physical_conditions}

% Because the source corpus contains objects with substantially different
% shapes and orientations, the fixture and loading regions are constructed
% automatically from the CAD geometry rather than manually specified for
% individual samples.

% For each solid, we first compute the centroid of every CAD surface. The
% pair of surface centroids with the largest Euclidean separation defines
% the principal analysis direction, which is normalized to obtain the
% functional analysis-axis unit vector. All surface centroids are then
% projected onto this axis. The projected extent of the solid is used to identify its two terminal
% surface groups. Surfaces whose centroid projections lie within \(5\%\)
% of the negative or positive end of the projected range are assigned to
% the corresponding terminal group. This surface-level grouping is used
% instead of selecting a single face, allowing geometries with multiple
% coplanar or fragmented terminal surfaces to be handled consistently.

% The negative terminal group is used as the fixture, with all
% translational degrees of freedom constrained,
% \(u_x=u_y=u_z=0\). The positive terminal group is used as the loading
% region. A uniform traction is distributed over this group such that its
% integrated resultant force has magnitude \(100\,\mathrm{kN}\) and is
% directed along the positive analysis axis.

% Both the analysis-axis vector and the resulting fixture and loading
% specifications are included in the task prompt. Therefore, although the
% boundary conditions are automatically derived from the reference
% geometry during dataset construction, they become explicit physical
% conditions that must also be satisfied by the generated design.


% \subsection{Inverse-Design Task Construction}
% \label{app:inverse_task}

% After a reference design has passed CAD construction and FEM analysis,
% we convert it into an inverse-design task. The conversion preserves the
% design family and physical requirements while deliberately removing the
% exact reference dimensions and simulation measurements.

% \paragraph{Coarse design domain.}
% We first compute the reference solid's global bounding-box extents
% \(e_x\), \(e_y\), and \(e_z\). Instead of exposing these values directly,
% each extent \(e\) is replaced by the enclosing interval formed by two
% adjacent powers of two:
% \begin{equation}
%     \mathcal{B}(e)
%     =
%     \left[
%         2^{\lfloor\log_2 e\rfloor},
%         2^{\lfloor\log_2 e\rfloor+1}
%     \right].
%     \label{eq:dataset_coarse_bound}
% \end{equation}
% The prompt therefore specifies only an admissible range for the global
% extent along each axis. For example, a reference extent lying between
% \(0.25\,\mathrm{m}\) and \(0.5\,\mathrm{m}\) is represented only by that
% interval rather than by its exact value.

% \paragraph{Performance requirements.}
% The displacement, stress, and mass requirements are derived from the
% corresponding reference measurements. For each measured quantity \(v\),
% we apply a \(10\%\) feasibility margin and round the resulting value
% upward to two significant digits:
% \begin{equation}
%     \tau(v)
%     =
%     \operatorname{RoundUp}_{2\mathrm{sig}}(1.10\,v).
%     \label{eq:dataset_threshold}
% \end{equation}
% The same procedure is applied independently to
% \(\delta_{\max}\), \(\sigma_{\mathrm{P99}}\), and \(m\).
% The resulting values become upper bounds on maximum displacement,
% \(99\)-th percentile von Mises stress, and total mass, respectively.

% The margin keeps the original reference design feasible under small
% numerical variations introduced by CAD processing, meshing, and FEM
% evaluation. Upward rounding additionally prevents the prompt from
% revealing the exact measurements of the reference design. 

% \subsection{Geometry Semantic Construction}
% \label{app:geometry_semantics}

% A central objective of the task construction is to preserve the intended
% geometric structure while preventing the model from directly recovering
% the original design parameters. We therefore construct the model-visible
% geometry semantics jointly from the source textual description and the
% CadQuery operation sequence.

% \paragraph{Numerical redaction.}
% We first sanitize the original textual description by removing numerical
% information that can reveal the reference geometry. This includes
% explicit dimensions, angles, ratios, percentages, quantities with units,
% and spelled-out numerical values. In contrast, non-numerical information
% describing the object category, shape composition, topology, feature
% types, feature counts, and spatial relationships is retained.

% For example, a source description may specify the exact diameter and
% height of a cylindrical feature. The sanitized description retains that
% a cylindrical feature exists and how it is positioned relative to the
% remaining geometry, but removes the original diameter and height. The
% generated program must determine such dimensions from the allowed design
% domain and physical requirements.

% \paragraph{Operation-derived semantics.}
% Textual descriptions alone do not always provide a sufficiently
% consistent representation of the underlying CAD construction. We
% therefore additionally inspect the CadQuery program and convert its
% modeling operations into dimension-free feature descriptions using a
% fixed operation-to-feature mapping. Numerical arguments of these
% operations are discarded.

% Table~\ref{tab:semantic_mapping} shows representative mappings used
% during this process. The mapping describes the geometric role of an
% operation rather than its parameters. For instance, a
% \texttt{cylinder} operation indicates the presence of a cylindrical
% feature, while a \texttt{cut} operation indicates a material-removal
% relationship. Sketches followed by extrusion retain the corresponding
% profile type and extrusion relation without preserving the original
% profile dimensions or extrusion distance.

% \begin{table}[t]
%     \centering
%     \small
%     \caption{
%     Representative CadQuery operation-to-feature mappings used to
%     construct dimension-free geometry semantics. Numerical parameters
%     associated with these operations are not exposed to the model.
%     }
%     \label{tab:semantic_mapping}
%     \setlength{\tabcolsep}{8pt}
%     \renewcommand{\arraystretch}{1.12}
%     \begin{tabular}{ll}
%         \toprule
%         \textbf{CadQuery construction} &
%         \textbf{Extracted geometry semantics} \\
%         \midrule
%         \texttt{box}
%         & rectangular-prismatic feature \\
%         \texttt{cylinder}
%         & cylindrical feature \\
%         polygonal sketch + \texttt{extrude}
%         & polygonal-prismatic feature \\
%         curved sketch + \texttt{extrude}
%         & curved-profile feature \\
%         \texttt{revolve}
%         & revolved feature \\
%         \texttt{cut}
%         & material-removal relation \\
%         \texttt{hole}
%         & hole or through-feature \\
%         \texttt{fillet}
%         & rounded-edge feature \\
%         \texttt{shell}
%         & hollow or shell construction \\
%         union / composition operations
%         & part-composition relation \\
%         \bottomrule
%     \end{tabular}
% \end{table}

% The sanitized source description and operation-derived information are
% then combined into a geometry-semantic description. The resulting
% description preserves the overall shape class, topology, feature types
% and counts, and relationships between individual components. Exact
% dimensions remain unspecified.

% For example, a reference composed of a hexagonal prism and a centered
% cylindrical top is described as requiring these two feature types, their
% relative placement, and their combination into a single solid, while the
% side lengths, cylinder diameter, and component heights remain hidden.
% Similarly, a design containing an extruded star-shaped profile combined
% with a cylindrical component preserves the corresponding feature
% composition without exposing the dimensions of either component.

% This construction makes the original reference a feasible realization
% of the requested design semantics, but not the uniquely specified
% solution. The model must select a concrete geometric realization that
% preserves the required structure while satisfying the design-domain and
% physical constraints.


% \subsection{Model-Visible and Private Information}
% \label{app:dataset_visibility}

% We maintain a strict separation between information used internally to
% construct each sample and information made available to the model. The
% model-visible prompt contains the dimension-free geometry semantics,
% coarse global bounding-box ranges, functional analysis axis, material
% properties, fixture and loading conditions, and the resulting
% displacement, stress, and mass limits.

% The exact reference dimensions are never included in the prompt.
% Likewise, the original CadQuery program, reference STEP model, FEM mesh,
% nodal displacement field, element stress field, and unrounded reference
% measurements remain private during generation. The reference CAD and FEM
% results therefore serve only as internal data for task construction and
% verification.

% This separation is particularly important for the inverse-design
% setting. The generated design is not evaluated by geometric similarity
% to the reference model. Instead, it must preserve the requested
% geometric semantics, remain inside the prescribed design domain, and
% satisfy the physical requirements through its own geometric realization.


% \subsection{Dataset Filtering and Partitioning}
% \label{app:dataset_split}

% After CAD execution, geometry processing, meshing, and FEM annotation,
% \(77{,}067\) source designs successfully pass the complete dataset
% construction pipeline.

% We partition these samples by sample ID into disjoint SFT and GRPO
% subsets. The SFT subset contains \(11{,}024\) samples. Each SFT sample
% pairs the constructed task prompt with its sanitized reference CadQuery
% program, allowing the model to first learn executable program
% construction and the required output format.

% The remaining \(66{,}043\) samples form the GRPO subset. Unlike SFT,
% these samples do not expose a target CadQuery program during
% reinforcement learning. Candidate programs are instead evaluated by the
% deterministic CAD--CAE verifier according to their executable validity,
% geometric compliance, physical feasibility, and design objective.

% Within the GRPO subset, \(62{,}670\) samples are used for reinforcement
% learning and \(3{,}373\) are reserved for evaluation. The SFT and GRPO
% sets contain no overlapping sample IDs. Importantly, the same task
% construction procedure and prompt format are used across the two stages;
% the difference lies only in whether a reference program is supplied as
% supervision or the generated program is evaluated through the verifier.


% \subsection{Representative Examples}
% \label{app:dataset_examples}

% Figure~\ref{fig:dataset_examples} presents representative samples
% generated by the complete pipeline. For illustration, each example
% shows the private reference geometry and its corresponding FEM von Mises
% stress field together with the final model-visible task prompt.

% The examples illustrate the distinction between the private reference
% design and the public inverse-design specification. High-level geometric
% structure is retained in the prompt, while exact dimensions are replaced
% by coarse admissible ranges. Similarly, the exact reference simulation
% results remain hidden and are converted into displacement, stress, and
% mass requirements. The model must therefore determine a new concrete
% set of dimensions through generation rather than reproduce the reference
% geometry directly.

\section{Implementation of Dataset Generation}
\label{app:dataset}

This appendix provides implementation details for the dataset-generation
procedure introduced in Sec.~\ref{sec:dataset}. We focus on the construction
of physical conditions, inverse-design constraints, and dimension-free
geometry semantics. 

\subsection{Reference CAD--CAE Annotation and Physical Conditions}
\label{app:dataset_annotation}

Each source CadQuery program is executed and retained only if it produces
exactly one valid connected solid that can be processed by the downstream
CAD--CAE pipeline. We analyze all retained designs using three-dimensional
small-strain linear elasticity in SI units, with isotropic structural steel
($E=2.1\times10^{11}$ Pa, $\nu=0.3$, and
$\rho=7850~\mathrm{kg/m^3}$). For each successfully simulated design, we
record the maximum nodal displacement $\delta_{\max}$, the volume-weighted
99-th percentile element von Mises stress $\sigma_{\mathrm{P99}}$, and the
total mass $m$.

Because the source geometries vary substantially in shape and orientation,
the fixture and loading regions are determined automatically. We compute the
centroid of every CAD surface and select the pair of surface centroids with
the largest Euclidean separation. The normalized direction between this pair
defines the functional analysis axis. All surface centroids are projected
onto this axis, and surfaces whose projections lie within $5\%$ of either end
of the projected range form the two terminal surface groups. This grouping
allows multiple coplanar or fragmented terminal surfaces to be handled
consistently.

The negative terminal group is fully fixed
($u_x=u_y=u_z=0$), while a uniform traction is applied to the positive
terminal group such that its integrated resultant force is $100$ kN along
the positive analysis axis. The resulting analysis axis, fixture definition,
and loading condition are explicitly included in the model prompt.

\subsection{Inverse-Design Task Construction}
\label{app:inverse_task}

Each successfully analyzed reference is converted into an inverse-design task
while withholding its exact dimensions and simulation measurements.

\paragraph{Coarse design domain.}
For each global bounding-box extent $e$, we expose only the enclosing interval
between adjacent powers of two,
\begin{equation}
    B(e)
    =
    \left[
    2^{\lfloor \log_2 e \rfloor},
    2^{\lfloor \log_2 e \rfloor+1}
    \right].
    \label{eq:bbox_domain}
\end{equation}
The model therefore receives a coarse admissible range along each axis rather
than the dimensions of the reference design.

\paragraph{Performance requirements.}
The displacement, stress, and mass limits are derived independently from
$\delta_{\max}$, $\sigma_{\mathrm{P99}}$, and $m$. Each reference measurement
is multiplied by $1.10$ and rounded upward to two significant digits to form
the corresponding upper bound. The margin keeps the reference design feasible
under small numerical variations from CAD processing, meshing, and FEM
evaluation, while rounding avoids exposing its exact simulation measurements.




\begin{figure*}[t]
    \centering
    \includegraphics[width=\textwidth]{fig/dataset_cad_fulu.pdf}
    \caption{
    \textbf{Representative physics-aware Text-to-CAD tasks.}
    Each example shows the private reference geometry and its FEM stress field
    together with the corresponding model-visible task. Exact reference
    dimensions and simulation measurements are withheld; the model receives
    dimension-free geometry semantics, coarse design bounds, physical
    conditions, and displacement, stress, and mass limits. It must generate
    an executable CadQuery design satisfying these requirements while
    minimizing mass.
    }
    \label{fig:dataset_gallery}
\end{figure*}




\subsection{Geometry Semantic Construction}
\label{app:geometry_semantics}

The geometry specification must preserve the intended design structure without
revealing the parameters of the reference CAD. We construct this specification
from both the source textual description and the CadQuery operation sequence.

\paragraph{Numerical redaction.}
We remove numerical information that may reveal the reference geometry,
including explicit dimensions, angles, quantities with
units, and spelled-out numerical values. We retain non-numerical information
describing the shape class, topology, feature types and counts, and spatial
or compositional relationships. Thus, for example, the existence and relative
placement of a cylindrical feature are preserved while its diameter and height
are withheld.

\paragraph{Operation-derived semantics.}
Because textual descriptions do not always capture the CAD construction
consistently, we additionally parse the CadQuery program and map its modeling
operations to dimension-free geometric features. Numerical operation arguments
are discarded. Representative mappings are shown in
Table~\ref{tab:operation_semantics}.

\begin{table}[t]
    \centering
    \caption{
    Representative CadQuery operation-to-feature mappings used to construct
    dimension-free geometry semantics. Numerical parameters are not exposed
    to the model.
    }
    \label{tab:operation_semantics}
    \small
    \begin{tabular}{ll}
        \toprule
        \textbf{CadQuery construction} & \textbf{Extracted geometry semantics} \\
        \midrule
        box                         & rectangular-prismatic feature \\
        cylinder                    & cylindrical feature \\
        polygonal sketch + extrude  & polygonal-prismatic feature \\
        curved sketch + extrude     & curved-profile feature \\
        revolve                     & revolved feature \\
        cut                         & material-removal relation \\
        hole                        & hole or through-feature \\
        fillet                      & rounded-edge feature \\
        shell                       & hollow or shell construction \\
        union / composition         & part-composition relation \\
        \bottomrule
    \end{tabular}
\end{table}

The sanitized textual description and operation-derived features are combined
into the final geometry-semantic specification. It preserves the overall shape
class, topology, feature composition, and relationships between components,
while leaving their exact dimensions unspecified. Consequently, the reference
CAD remains one feasible realization of the task rather than a uniquely
specified target geometry.

Representative tasks produced by the complete pipeline are shown in
Fig.~\ref{fig:dataset_gallery}, where the reference CAD and FEM results are
included only for illustration and are not provided to the model.

% \begin{figure*}[t]
%     \centering
%     \includegraphics[width=\textwidth]{fig/dataset_cad_fulu.pdf}
%     \caption{
%     \textbf{Representative physics-aware Text-to-CAD tasks.}
%     Each example shows the private reference geometry and its FEM von
%     Mises stress field together with the corresponding model-visible
%     task prompt. Exact reference dimensions and simulation measurements
%     are removed: the geometry is represented by dimension-free semantic
%     descriptions and coarse bounding-box ranges, while the reference
%     physical measurements are converted into displacement, stress, and
%     mass limits. The model must generate an executable CadQuery program
%     that preserves the specified geometric structure, satisfies all
%     physical requirements, and minimizes mass. The reference geometry
%     and FEM visualization are shown only for illustration and are not
%     provided to the model.
%     }
%     \label{fig:dataset_examples}
% \end{figure*}




\section{Hierarchical Balanced GRPO}
\label{app:hierarchical_grpo}

Algorithm~\ref{alg:hierarchical_grpo} summarizes the training procedure of our
Hierarchical Balanced GRPO. 

\begin{algorithm}[h]
\caption{\textbf{Hierarchical Balanced GRPO.}}
\label{alg:hierarchical_grpo}
\begin{algorithmic}[1]
\REQUIRE Policy $\pi_\theta$, reference policy $\pi_{\mathrm{ref}}$,
prompts $\mathcal{Q}$, group size $G$, weights $\lambda_L,\lambda_w,\beta$

\FOR{each training iteration}
    \STATE Sample a prompt batch $\mathcal{B}\subset\mathcal{Q}$

    \FOR{each prompt $q\in\mathcal{B}$}
        \STATE Sample $\{o_i\}_{i=1}^{G}\sim\pi_\theta(\cdot\mid q)$

        \STATE \textbf{Verification:}
        evaluate each $o_i$ with the CAD--CAE verifier to obtain
        level $L_i\in\{0,1,2,3\}$ and level-specific scores
        $\{s_{ik}\}_{k=1}^{K_{L_i}}$
        \STATE \hspace{\algorithmicindent}
        $L_i=0$: invalid CAD;\;
        $L_i=1$: geometry failure;\;
        $L_i=2$: physics failure;\;
        $L_i=3$: fully feasible

        \STATE \textbf{Between-level advantage:}
        $\bar L=\frac{1}{G}\sum_{j=1}^{G}L_j$,\quad
        $A_i^{\mathrm{level}}=L_i-\bar L$

        \STATE \textbf{Within-level PeerRank:}
        initialize $\widetilde A_i=0$ for all $i$
        \FOR{$\ell=0,1,2,3$}
            \STATE $I_\ell=\{i:L_i=\ell\}$
            \IF{$|I_\ell|>1$}
                \STATE Compute tie-aware centered ranks
                $r_{ik}$ of $\{s_{ik}\}_{i\in I_\ell}$
                for $k=1,\ldots,K_\ell$
                \STATE $\widetilde A_i
                =\sum_{k=1}^{K_\ell}w_{\ell k}r_{ik}$,
                \quad $\forall i\in I_\ell$
            \ENDIF
        \ENDFOR

        \STATE \textbf{Hierarchical advantage:}\quad
        $A_i=\lambda_L A_i^{\mathrm{level}}
        +\lambda_w\widetilde A_i$
    \ENDFOR

    \STATE Update $\theta$ using DAPO-style GRPO:
    $\mathcal{L}
    =\mathcal{L}_{\mathrm{DAPO}}(\theta;A)
    +\beta D_{\mathrm{KL}}(\pi_\theta\|\pi_{\mathrm{ref}})$
\ENDFOR

\RETURN Optimized policy $\pi_\theta$
\end{algorithmic}
\end{algorithm}
